\documentclass[10pt,conference]{IEEEtran}
\IEEEoverridecommandlockouts
% The preceding line is only needed to identify funding in the first footnote. If that is unneeded, please comment it out.
\usepackage{cite}
\usepackage{amsmath,amssymb,amsfonts}
\usepackage{algorithmic}
\usepackage{graphicx}
\usepackage{textcomp}
\usepackage{xcolor}
\usepackage{url}
% \usepackage{hyperref}
% \usepackage{svg}
\usepackage[inline]{enumitem}
\usepackage{booktabs}

% custom packages
\usepackage{subfiles}
\usepackage{physics}
\usepackage{listings}
\usepackage{multirow}
\usepackage{longtable}
\usepackage{comment}
\usepackage{tabularx}
\usepackage[compatibility=false]{caption}
\usepackage{capt-of}
\usepackage{orcidlink}

% Define custom colors for Python syntax highlighting
\definecolor{codegreen}{rgb}{0,0.6,0}       % Green for comments
\definecolor{codegray}{rgb}{0.5,0.5,0.5}        % Gray for line numbers
\definecolor{codepurple}{rgb}{0.58,0,0.82}     % Purple for strings and special literals
\definecolor{backcolour}{rgb}{0.95,0.95,0.92}   % Light beige for background
\definecolor{builtincolor}{rgb}{0.2, 0.2, 0.7}   % Blue for built-ins
\definecolor{defclasscolor}{rgb}{0.1, 0.5, 0.1} % Dark Green for def/class
\definecolor{docstringcolor}{rgb}{0.25,0.5,0.35} % Desaturated green for docstrings
% Standard magenta (LaTeX built-in) is used for general keywords

\lstdefinestyle{pythonStyle}{
    language=Python,
    backgroundcolor=\color{backcolour},
    basicstyle=\ttfamily\footnotesize,
    commentstyle=\color{codegreen},
    stringstyle=\color{codepurple},
    keywordstyle=\color{magenta},
    numbers=left,
    numberstyle=\tiny\color{codegray},
    numbersep=5pt,
    xleftmargin=\parindent,
    tabsize=2,
    breaklines=true,
    breakatwhitespace=false,
    showspaces=false,
    showstringspaces=false,
    showtabs=false,
    keepspaces=true,
    captionpos=bl
}

\lstset{style=pythonStyle}

\def\BibTeX{{\rm B\kern-.05em{\sc i\kern-.025em b}\kern-.08em
    T\kern-.1667em\lower.7ex\hbox{E}\kern-.125emX}}

\hyphenation{Hellinger Bhattacharyya Michelson Estimator Sampler}

\newcommand{\TVD}{\mathrm{TVD}}
\newcommand{\BC}{\mathrm{BC}}
\newcommand{\CCS}{\mathrm{CCS}}
\newcommand{\HF}{F_{H}}

\newtheorem{lemma}{Lemma}
\newtheorem{proposition}{Proposition}

\begin{document}

\title{A Representational Analysis of Quantum Execution Success Measures}

% \author{\IEEEauthorblockN{1\textsuperscript{st} Cristian Marquez}
% \IEEEauthorblockA{\textit{Department of Systems and Computing Engineering} \\
% \textit{Universidad de los Andes}\\
% Bogotá, Colombia \\
% c.marquezb$@$uniandes.edu.co}
% \and
% \IEEEauthorblockN{2\textsuperscript{st} Daniel Sierra}
% \IEEEauthorblockA{\textit{Electrical Engineering and Computer Science} \\
% \textit{The Catholic University of America, Washington}\\
% DC, 20064 \\
% sierrasosa$@$cua.edu}
% \and
% \IEEEauthorblockN{3\textsuperscript{nd} Kelly Garcés}
% \IEEEauthorblockA{\textit{Department of Systems and Computing Engineering} \\
% \textit{Universidad de los Andes}\\
% Bogotá, Colombia \\
% kj.garces971$@$uniandes.edu.co}
% }

\author{\IEEEauthorblockN{1\textsuperscript{st} Cristian Marquez \orcidlink{0009-0003-2708-0469}}
\IEEEauthorblockA{\textit{Department of Systems} \\
\textit{and Computing Engineering}\\
\textit{Universidad de los Andes}\\
Bogotá, Colombia \\
ing.cristian.marquez$@$gmail.com}
\and
\IEEEauthorblockN{2\textsuperscript{nd} Daniel Sierra \orcidlink{0000-0003-1326-0867}}
\IEEEauthorblockA{\textit{Computer Science} \\
\textit{The Catholic University of America}\\
Washington DC, 20064 \\
sierrasosa$@$cua.edu}
\and
\IEEEauthorblockN{3\textsuperscript{rd} Kelly Garc\'es \orcidlink{0000-0002-1070-439X}}
\IEEEauthorblockA{\textit{Department of Systems} \\
\textit{and Computing Engineering}\\
\textit{Universidad de los Andes}\\
Bogotá, Colombia \\
kj.garces971$@$uniandes.edu.co}
}


\maketitle

\begin{abstract}

Evaluating quantum execution success requires identifying the behavior expected from the program. Quantities such as success rate, Hellinger fidelity, and total variation distance can share a numerical range while representing different attributes and supporting different comparisons. We apply representational measurement theory and the Goal--Question--Metric approach to classify quantities associated with Sampler and Estimator results by their attributes, reference requirements, and permitted operations.

For a fixed attribute and reference, strictly monotone transformations preserve individual threshold decisions when thresholds and inequality directions are converted consistently. Nonlinear aggregation can nevertheless reverse a group ranking. We further show that coarse Hellinger fidelity combines accepted answer probability with agreement about the relative frequencies of accepted answers, while differential success rate represents a censored ordering of peak dominance. These quantities therefore need not support the same conclusion about task success.

We examine these distinctions using a corpus of 1,478 execution records, including hardware and simulation results, and a hardware Estimator execution. Stored distance and similarity pairs reproduce their defining transformations, and no aggregation reversal occurs in 31 selected provider comparisons. Grover cases produce different decisions under task and complete ideal references, while the Estimator record exposes numerical reports whose labels do not represent the intended agreement. These findings support a procedure for selecting and reporting execution quantities according to the program requirement, available reference, and intended operation.

\end{abstract}

\begin{IEEEkeywords}
Quantum software measurement, representational measurement theory, Goal--Question--Metric, Hellinger fidelity, total variation distance, Sampler, Estimator
\end{IEEEkeywords}

\section{Introduction}
\label{sec:intro}

The result of a quantum processing unit (QPU) does not determine by itself what success means. In current quantum execution frameworks, its interpretation depends on the measurements performed and the Primitive through which their results are reported. A Sampler Job stores measurement results as samples from output registers, while an Estimator Job aggregates results into expectation values for chosen observables~\cite{qiskit2024}, as described in Section~\ref{sec:primitives}. Because the two outputs preserve different empirical relations, applying a formula defined for one interpretation to evaluate the other changes the attribute being measured, not only the notation~\cite{fentonBieman2015}.

Quantum information already provides a distinction between fidelity and distance, along with metrics computed using multiple fidelity functions~\cite{gilchrist2005,luo2004,ma2009,nielsen2010}. The unresolved problem in software measurement is how to select among success rate, Hellinger fidelity (HF), total variation distance (TVD), probability of successful trials (PST), estimated success probability (ESP) and related scores when analyzing the output of a quantum Job. These quantities are commonly averaged, thresholded, or compared even when they do not share a domain.

For a Job evaluated against a fixed reference, the conditions $\TVD\le0.2$ and $1-\TVD\ge0.8$ yield the same threshold decision. The second quantity is a decreasing transformation of the first, and both the threshold and inequality direction are converted consistently. Section~\ref{sec:threshold} formalizes this equivalence.

Equivalent decisions for individual Jobs do not guarantee equivalent conclusions under other operations. Averaging quantities related by a nonlinear transformation can change a group ranking. Propagating error bounds, interpreting ratios, and applying thresholds across tasks also require justification based on the represented attribute and the operation performed. When the required inputs or reference information are absent, the corresponding comparison should be reported as unavailable. Selecting an execution measure therefore requires establishing both what it represents and which operations preserve the intended conclusion.

This paper classifies quantities used to evaluate a Job through representational measurement theory and the Goal--Question--Metric (GQM) approach. It establishes when transformations preserve execution conclusions and when subsequent operations require additional justification. The analysis also characterizes differential success rate (DSR) as an ordinal score of peak dominance and examines these distinctions using QPU execution data.

The study addresses three research questions (RQs).
\begin{enumerate}[label=\textbf{RQ\arabic*},ref=\arabic*,leftmargin=*,itemsep=2pt,parsep=0pt,topsep=3pt]
\item\label{rq:attributes} Which execution attribute and reference information does each quantity represent?
\item\label{rq:operations} Which transformations and operations preserve a conclusion, and which can change it?
\item\label{rq:evidence} Do these distinctions appear in QPU execution data, and how do they inform the interpretation of DSR?
\end{enumerate}

The analysis separates two questions that are easily confused. The mathematical question is whether a function satisfies Axioms M1--M4 on a named space. The representational question asks whether a numerical value preserves the intended relations for an attribute of a specified entity~\cite{fentonBieman2015,fenton1994}. Passing one test does not guarantee passing the other. For example, Hellinger distance is a metric on distributions, but agreement between histograms may be the wrong attribute when the decision concerns only the probability mass assigned to an answer set $E$.

The remainder of this paper is organized as follows. Section~\ref{sec:background} reviews the theoretical basis, and Section~\ref{sec:related} positions the work against related studies. Section~\ref{sec:model} classifies execution quantities by their attributes and reference requirements. Section~\ref{sec:analysis} establishes the effects of transformations and operations on their interpretation, and Section~\ref{sec:evaluation} examines these distinctions in execution data. Sections~\ref{sec:discussion}--\ref{sec:conclusion} provide reporting guidance, limitations, and conclusions.

\section{Background}
\label{sec:background}

The quantities compared in this study can share a numerical range while following different mathematical rules. In this section, we first establish their domains and properties and then connect those properties to the execution attributes they represent.

\subsection{Distance Metrics and Related Quantities}
\label{sec:vocab}

Sampler outcomes and ideal references for the same fixed measurement and outcome alphabet can be represented as probability distributions over a finite set of outcomes. For $n$ possible outcomes, their common domain is the probability simplex

\begin{equation}
\Delta_n=
\left\{P\in\mathbb{R}^{n}\mid
P(x)\ge0,\ \sum_{x=1}^{n}P(x)=1\right\}.
\label{eq:probability-simplex}
\end{equation}

The value $P(x)$ is the probability assigned to the outcome $x$. The Bhattacharyya coefficient (BC)~\cite{bhattacharyya1943} measures the overlap between two distributions in this simplex:

\begin{equation}
\BC(P,Q)=\sum_{x=1}^{n}\sqrt{P(x)Q(x)}.
\label{eq:bhattacharyya-coefficient}
\end{equation}

The classical Hellinger distance and TVD are

\begin{equation}
H(P,Q)=\sqrt{1-\sum_{x=1}^{n}\sqrt{P(x)Q(x)}}
=\sqrt{1-\BC(P,Q)},
\label{eq:hellinger-distance}
\end{equation}
\begin{equation}
\TVD(P,Q)=\tfrac12\sum_{x=1}^{n}\lvert P(x)-Q(x)\rvert.
\label{eq:tvd}
\end{equation}

The HF used in this paper is the squared overlap

\begin{equation}
\HF(P,Q)=\BC(P,Q)^{2}=(1-H(P,Q)^{2})^{2},
\label{eq:hellinger-fidelity}
\end{equation}

Equation~\ref{eq:hellinger-fidelity} shows that HF is a function of $H^2$, not the complement $1-H$~\cite{cha2007}, whereas TVD similarity is the complement of TVD.

\begin{equation}
\sigma_{\TVD}(P,Q)=1-\TVD(P,Q)
\label{eq:tvd-similarity-definition}
\end{equation}

Their norm representations are

\begin{align}
H(P,Q)&=\tfrac{1}{\sqrt{2}}\lVert\sqrt{P}-\sqrt{Q}\rVert_{2},
\label{eq:hellinger-l2}\\
\TVD(P,Q)&=\tfrac{1}{2}\lVert P-Q\rVert_{1}.
\label{eq:tvd-l1}
\end{align}

These norm representations establish that both $H$ and TVD are metrics on $\Delta_n$~\cite{gilchrist2005,nielsen2010}. More generally, on a set $X$, a distance metric is a function
\begin{equation}
 d\colon X\times X\rightarrow\mathbb{R}
 \label{eq:metric-map}
\end{equation}
that satisfies the following four axioms:

\begin{align}
 d(x,y) &\ge 0. &\tag{M1}\label{eq:metric-nonnegativity}\\
 d(x,y) &= 0\iff x=y. &\tag{M2}\label{eq:metric-identity}\\
 d(x,y) &= d(y,x). &\tag{M3}\label{eq:metric-symmetry}\\
 d(x,z) &\le d(x,y)+d(y,z). &\tag{M4}\label{eq:metric-triangle}
\end{align}

With Axioms~M1--M4 in place, the quantities commonly reported as quantum metrics can be compared with the formal definition of a distance metric. Among the quantities considered in this paper, Hellinger distance and TVD satisfy M1--M4 on $\Delta_n$ through their norm representations. Other notions appearing in the literature have different mathematical properties. A \emph{pseudometric} relaxes the identity of indiscernibles (M2) by allowing distinct points to have zero separation, while a \emph{quasimetric} relaxes symmetry (M3). A divergence such as Kullback--Leibler is typically asymmetric and may violate the triangle inequality (M4).

By contrast, HF and TVD similarity---the similarity quantities used in our examples---are maximized when their two arguments coincide. Other commonly reported quantities still have different forms: the success rate and the PST are proportions, the ESP is a prediction, and DSR is evaluated in this paper as an ordinal score. These quantities evaluate a Job relative to a goal, reference, or calibration model and are therefore not candidates for Axioms~M1--M4 in their original form. This distinction has two practical consequences. First, taking the complement of a distance does not automatically produce another metric. Second, any claim that invokes the triangle inequality should be stated for a verified distance metric or for a dissimilarity explicitly induced from one. For example, the Hellinger construction in Equation~\ref{eq:hellinger-distance} converts the Bhattacharyya overlap of Equation~\ref{eq:bhattacharyya-coefficient} into a metric. Another standard conversion is the angular distance
\begin{equation}
\theta(P,Q)=\arccos\BC(P,Q).
\label{eq:bhattacharyya-angle}
\end{equation}
Both constructions are discussed in~\cite{luo2004,ma2009}. Accordingly, we use \emph{distance metric} only for a function satisfying Axioms~M1--M4 on a named space, and use \emph{measure}, \emph{score}, or \emph{similarity} only after naming the attribute.

The preceding definitions are for classical outcome distributions. For comparisons at the state level, we use the squared Uhlmann fidelity convention
\begin{equation}
F(\rho,\sigma)=
\left(\mathrm{Tr}\sqrt{\sqrt{\rho}\sigma\sqrt{\rho}}\right)^2.
\label{eq:uhlmann-fidelity}
\end{equation}
In what follows, we explicitly distinguish similarities from distances and name the underlying space whenever invoking metric properties such as the triangle inequality. This convention makes the mathematical meaning of \emph{closeness} explicit before its suitability as a representation of an execution attribute is assessed.

\subsection{Representational Measurement Theory}
\label{sec:fenton}

The preceding subsection establishes whether a quantity is a mathematical metric in a named space. In this section, we establish whether a numerical value preserves the intended relations for a specified execution attribute. Fenton and Bieman~\cite{fentonBieman2015} define measurement as a mapping from an empirical relational system to a numerical one that preserves those relations, a requirement known as the representation condition~\cite{polancic2017complexity}. Applied to a quantum Job, this view requires naming the attribute being measured, such as proximity to a reference distribution, probability mass on $E$, dominance of expected peaks, or deviation from an ideal expectation.

Table~\ref{tab:scale-types} summarizes the five types of scale used in this paper and the comparisons supported by each.

\begin{table}[ht]
\centering
\caption{Measurement scale types and their interpretations.}
\label{tab:scale-types}
\renewcommand{\arraystretch}{1.08}
\begin{tabularx}{\columnwidth}{@{}p{0.16\columnwidth}X@{}}
\toprule
\textbf{Scale} & \textbf{Interpretation} \\
\midrule
Nominal & Labels only. A provider name is not a success quantity. \\
Ordinal & Order is preserved; differences are not. Against the same reference and on a common outcome space, a ranking of Jobs by $\HF$ supports "\textit{closer than,}" not "\textit{twice as close.}" \\
Interval & Differences are meaningful, whereas the zero point is conventional. Treating an expectation value as measured on an interval scale requires an explicit empirical interpretation. \\
Ratio & A meaningful, nonarbitrary zero and admissible transformations restricted to positive scaling are required to support statements such as "\textit{twice as far.}" Satisfaction of Axioms M1--M4 alone does not establish a ratio scale, and $1-d$ does not inherit one automatically. \\
Absolute & Shot counts correspond most closely to this type. Derived proportions inherit meaning only after the derivation is stated~\cite{fentonBieman2015}. \\
\bottomrule
\end{tabularx}
\end{table}

A scale type belongs to a stated attribute and measurement procedure, not to the numerical value alone. Table~\ref{tab:audit} therefore assigns a conservative measurement status rather than inferring a scale from a familiar formula. Subsection~\ref{sec:operations} applies this principle to $p_E$, aggregation between Jobs, ratio statements, and uncertainty due to finite sampling. Mathematical kind, scale interpretation, and sampling uncertainty remain separate properties.

This framework prevents every numerical quantity from being treated as a mathematical metric, a concern raised by Fenton~\cite{fenton1994,fenton1992}. A unary score can represent the probability mass on an answer set, even though Axioms M1--M4 do not apply to it. On the other hand, a mathematical distance can be unsuitable for a decision concerning a different attribute. Clipping a score corrected for chance at zero provides another limitation because it removes the ordering among Jobs below the chance baseline.

Representational measurement theory determines the interpretations and comparisons that a quantity supports. The GQM paradigm complements this analysis by guiding the selection of a quantity from the execution goal and available information~\cite{basili1994,fentonBieman2015}. It also distinguishes measurements from prediction; for example, ESP predicts execution success from calibration data rather than measuring an observed histogram. Section~\ref{sec:related} contrasts this question with the evaluation of predictive validity. Section~\ref{sec:selection} applies GQM by naming the object, purpose, quality focus, viewpoint, and execution context before deriving questions and quantities.

\section{Related Work}
\label{sec:related}

The literature on quantum information establishes the distinctions of fidelity and distance~\cite{gilchrist2005,luo2004,ma2009,mendonca2008,uhlmann1976,nielsen2010}, and we summarize them in Section~\ref{sec:vocab}. Since these results already distinguish similarities from the distances they induce, this research therefore does not claim the metric axioms or the fidelity distinction as new results.

The most closely related study to ours is by Mohapatra \emph{et al.}~\cite{mohapatra2025}, who link HF, TVD, ESP, and PST to the fidelity of density matrices. Their strongest reported proxies---ESP and PST---answer how well an inexpensive quantity tracks a reference obtained through tomography. However, Mohapatra claims are complementary to, but distinct from, deciding which attribute should define success for a quantum Job.

Within quantum software engineering, multiple surveys describe distribution distances, statistical tests, assertions, and other criteria used to decide whether a quantum execution succeeded~\cite{paltenghi2024,andrews2026,ramalho2025}. Ye \emph{et al.}~\cite{ye2025} differentiate the validation of measured output values from the validation of complete distributions and show how finite sampling constrains both. The work of Ye establishes the testing context in which a success quantity is consumed, but does not provide the integrated classification framework developed in this study.

Hacaloglu \emph{et al.}~\cite{hacaloglu2026} apply the broader questions of what, how, and why to the measurement of quantum software size. That study targets an internal product attribute rather than execution success, but it demonstrates the relevance of classical software measurement theory to quantum computing.

In contrast, this paper focuses on quantities derived from quantum Job outcomes and provides both a classification framework and a procedure based on GQM for selecting a success quantity. Our literature review did not identify prior work that integrates these elements, so our contribution is primarily based on scope and integration rather than novelty claims for the underlying distance measures or axioms.

\section{Classification of Quantum Execution Quantities}
\label{sec:model}

Evaluating execution success requires identifying the behavior expected from the quantum program. A task may require an accepted classical answer, agreement with an ideal distribution, or an expectation value within a specified tolerance. These requirements concern different execution attributes, and the available Primitive result determines which of them can be evaluated directly.

In Qiskit Runtime, \emph{SamplerV2} returns samples from classical output registers, whereas \emph{EstimatorV2} returns expectation estimates and their uncertainties for specified observables~\cite{qiskit2024,ibmEstimator2026}. These views can involve different measurements and processing, including basis rotations and mitigation. Table~\ref{tab:primitives} connects each result to the reference needed to assess agreement.

\begin{table}[ht]
\centering
\caption{Questions for evaluating Sampler and Estimator results.}
\label{tab:primitives}
\renewcommand{\arraystretch}{1.15}
\begin{tabularx}{\columnwidth}{@{}>{\raggedright\arraybackslash}p{0.22\columnwidth}>{\raggedright\arraybackslash}X>{\raggedright\arraybackslash}X@{}}
\toprule
\textbf{Question} & \textbf{Sampler} & \textbf{Estimator} \\
\midrule
What result is available? & Samples in classical registers, summarized by $\hat{P}$
& Estimates $\widehat{\langle O_{i}\rangle}$ for specified observables \\
What reference is needed? & A goal set $E$ or a full ideal distribution $P^{\star}$
& Ideal expectation values $\langle O_i\rangle$ \\
What can be evaluated? & Did the measured outcomes agree with $E$ or $P^{\star}$?
& Did the estimated expectation values agree with $\langle O_i\rangle$? \\
\bottomrule
\end{tabularx}
\end{table}

\subsection{Interpreting Sampler and Estimator Results}
\label{sec:primitives}

When designing an execution, the evaluation goal guides the choice of measurements and Primitive. For a completed Job, the available result constrains which attributes can still be assessed. In either case, a quantity requires both a compatible result and the reference information appropriate to the goal, as summarized in Table~\ref{tab:primitives}.

For a Sampler result with $S>0$ shots, normalizing the counts $\mathcal C(x)$ gives the empirical distribution $\widehat P$. A set of accepted answers $E$ defines the observed task success $p_E$. A justified null distribution $Q_0$ additionally supports chance corrected success (CCS), provided its baseline $b<1$.
\begin{align}
\widehat P(x)&=\frac{\mathcal C(x)}{S},\qquad K=|E|, \notag\\
p_E&=\sum_{x\in E}\widehat P(x)
=\frac{1}{S}\sum_{x\in E}\mathcal C(x),
\label{eq:success-mass}\\
b&=\Pr_{X\sim Q_0}(X\in E),
\label{eq:chance-baseline}\\
\CCS&=\mathrm{clip}\left(\frac{p_E-b}{1-b},0,1\right).
\label{eq:chance-corrected-success}
\end{align}

Equation~\ref{eq:chance-corrected-success} measures success beyond the selected baseline and assigns zero to values at or below it. Uniform guessing over $m$ measured bits gives $b=K/2^m$. A complete ideal distribution $P^{\star}$ supports a different question about agreement across the full outcome space. Its probabilities are supplied by the ideal execution, rather than chosen from the observed counts.

A task can instead prescribe weights $w_i\geq0$ with $\sum_iw_i=1$ for its accepted answers $E=\{e_1,\ldots,e_K\}$. The task target $W_E$ assigns $w_i$ to $e_i$ and zero outside $E$. A coarse comparison retains each accepted answer and combines all other outcomes into one bin, comparing $(\widehat P(e_1),\ldots,\widehat P(e_K),1-p_E)$ with $(w_1,\ldots,w_K,0)$. Uniform weights express a requirement for balanced accepted outcomes, which is stronger than accepting any answer in $E$. Section~\ref{sec:threshold} establishes the consequence for task success.

An Estimator result evaluates agreement for specified observables. Given an ideal expectation, absolute deviation and the observable fidelity used here express this agreement as

\begin{equation}
d_i=\left|\widehat{\langle O_i\rangle}-\langle O_i\rangle^{\mathrm{ideal}}\right|,
\qquad
F_{\mathrm{obs},i}=1-\frac{d_i}{2}.
\label{eq:observable-fidelity}
\end{equation}

For ideal and estimated expectations in $[-1,1]$, $F_{\mathrm{obs},i}$ lies in $[0,1]$. Mitigation can move an estimate outside this interval, so the score is not always bounded. For an observable with eigenvalues in $\{\pm1\}$, an unmitigated estimate also gives the positive eigenspace mass through $p_+(O)=(1+\widehat{\langle O\rangle})/2$. This mass represents task success only when that eigenspace defines the goal. A result containing only \texttt{evs} and \texttt{stds} supplies neither the histogram needed for HF nor the competing peaks needed for DSR.

Table~\ref{tab:references} names the reference classes used throughout the analysis. They include operational protocols for comparisons that require additional executions or measurements.

\begin{table}[ht]
\centering
\caption{Reference classes used to interpret execution quantities.}
\label{tab:references}
\renewcommand{\arraystretch}{1.15}
\begin{tabularx}{\columnwidth}{@{}l>{\raggedright\arraybackslash}p{0.37\columnwidth}>{\raggedright\arraybackslash}X@{}}
\toprule
\textbf{Class} & \textbf{What information is needed?} & \textbf{What does it support?} \\
\midrule
R1 task & Answer set $E$, with weights $w$ or null model $Q_0$ when required
& Task success, correction for chance, or a coarse distribution comparison \\
R2 distribution & Complete ideal distribution $P^{\star}$
& Comparison of complete histograms \\
R3 observable & Ideal expectations, or a Hamiltonian and reference energy
& Agreement for observable values or energy \\
R4 operational & Fidelity protocol (R4F) or inverse circuit (R4I)
& State, process, or inverse return comparison \\
\bottomrule
\end{tabularx}
\end{table}

Some quantities use other information. ESP needs a calibration model, while the standardized magnitude $z_0=|\widehat{\langle O\rangle}|/s_O$ needs a positive reported uncertainty $s_O$. Neither establishes agreement with an ideal expectation. Relative error instead compares absolute deviation with the magnitude of a nonzero ideal. These requirements distinguish predictions and diagnostics from evaluations against R1 through R4.

% This subsection operationalizes RQ1, and the permitted use column informs RQ2.
\subsection{Attributes and Permitted Comparisons}
\label{sec:classification}

For each quantity, we identify its attribute and inputs, assess its mathematical kind on the named domain, and state the comparisons supported by that interpretation. Table~\ref{tab:audit} applies these criteria using the reference classes in Table~\ref{tab:references}. Its permitted uses follow the distinction between metric properties and measurement scales established in Section~\ref{sec:fenton}.

\begin{table*}[ht]
\centering
\caption{Representational classification of execution quantities. Bold entries satisfy Axioms M1--M4 on the named space. Other entries are assessed by the relation they preserve.}
\label{tab:audit}
\scriptsize
\renewcommand{\arraystretch}{1.12}
\begin{tabularx}{\textwidth}{@{}>{\raggedright\arraybackslash}p{0.15\textwidth}
>{\raggedright\arraybackslash}p{0.13\textwidth} c
>{\raggedright\arraybackslash}p{0.10\textwidth}
>{\raggedright\arraybackslash}p{0.22\textwidth}
>{\raggedright\arraybackslash}X@{}}
\toprule
\textbf{Quantity} & \textbf{Attribute} & \textbf{Ref.} & \textbf{Kind}
& \textbf{Status and permitted use} & \textbf{Principal limitation} \\
\midrule
Success rate $p_{E}$ & mass on $E$ & R1($E$) & proportion
& derived proportion; order and thresholds; state how Jobs or shots are averaged
& ignores competitors and variation within $E$ \\
CCS & mass beyond a null baseline & R1($E,Q_0$) & censored score
& order only above $b$; thresholds require the same null model
& clipping removes all order below $b$ \\
Hellinger distance $H$ & histogram distance & R2 & \textbf{metric}
& order; descriptive means require a stated estimand and weighting; triangle bounds after valid maps
& needs $P^{\star}$; worst case $O(2^{m})$ representation \\
HF $\HF$ & histogram overlap & R2 & similarity
& ordinal within a fixed reference; converted thresholds and rank summaries
& arithmetic means need not preserve distance rankings \\
TVD & $L^{1}$ histogram distance & R2 & \textbf{metric}
& order; descriptive means require a stated estimand and weighting; triangle bounds
& needs $P^{\star}$; ignores phase \\
TVD similarity & histogram similarity & R2 & similarity
& exact complement; converted thresholds and ranks
& metric chaining uses the underlying TVD \\
Uhlmann fidelity & state overlap & R4F & similarity
& overlap order under a fixed protocol
& not a bitstring score; protocol assumptions matter \\
Bures / sine distances & state distance & R4F & \textbf{metrics}
& distance order and triangle bounds on states
& inherit the fidelity protocol requirements \\
ESP & predicted hardware success & calibration & prediction
& evaluate calibration and predictive validity
& not an observed execution attribute \\
PST & return mass after $U^{\dagger}U$ & R4I & proportion
& order and thresholds for the inverse construction
& evaluates a different circuit; relaxation can favor zero \\
Coarse TVD & distance on $\Delta_{K+1}$ & R1($E,w$) & \textbf{metric}
& distance operations on the coarse space
& generally a pseudometric on full histograms; $w$ must be justified \\
Coarse HF & coarse overlap & R1($E,w$) & similarity
& ordinal; convert to coarse $H$ for distance operations
& reduces to $p_E$ at $K=1$ \\
DSR (Michelson) & peak dominance & R1($E$) & ordinal representation
& order for a fixed task and sampling design
& clipping; mean versus maximum; depends on $K$ \\
Absolute expectation distance & expectation difference & R3 & \textbf{metric on $\mathbb{R}$}
& absolute differences and triangle bounds for a fixed observable
& Hamiltonian aggregation uses coefficients and signed deviations \\
Observable fidelity & expectation similarity & R3 & similarity
& affine threshold conversion per observable
& aggregation requires stated observable weights \\
Positive eigenspace mass & mass for eigenvalue $+1$ & observable $O$ & proportion
& physical proportion for an unmitigated estimate of involutory $O$
& task success only when the positive eigenspace defines the goal \\
Relative error & relative expectation deviation & R3 & derived ratio
& compare only for nonzero ideals on a shared interpretation
& undefined and unstable at zero \\
Standardized magnitude $z_0$ & magnitude relative to zero & uncertainty & diagnostic
& standardized by the reported uncertainty, not task success
& does not alone establish significance or agreement \\
\bottomrule
\end{tabularx}
\end{table*}

The success rate and HF entries in Table~\ref{tab:audit} illustrate the consequence for program evaluation. Success rate concerns accepted answers, while fidelity concerns agreement with a distribution whose probabilities may encode additional requirements. A developer must therefore fix the intended attribute before comparing values. Even when two quantities represent the same attribute, the operations applied to them can affect the resulting conclusion.

\section{Effects of Transformations and Operations on Execution Conclusions}
\label{sec:analysis}

% Internal review guide. This section supplies the analytical answer to RQ2 by distinguishing transformations that preserve a conclusion from operations that require additional assumptions. The Evaluation tests selected consequences with execution data.
An execution accepted under a distance threshold should receive the same decision when that distance is expressed as an equivalent similarity and the threshold is converted consistently. This equivalence concerns a fixed attribute and reference. The following results establish its scope, distinguish task success from peak dominance, and examine what changes when values are aggregated or used to bound other comparisons.

\subsection{Equivalent Decisions for a Fixed Interpretation}
\label{sec:threshold}

A fixed attribute and reference, a strictly monotone transformation, and a consistently converted threshold are sufficient to preserve a decision for an individual Job.

\begin{lemma}[Threshold invariance]
\label{lem:threshold}
Let $d(\cdot,Q)$ be a real quantity with fixed reference $Q$, and let $\phi$ be strictly monotone on an interval containing its values and the threshold $\tau$. Then
\begin{equation}
d(P,Q)\leq\tau
\iff
\begin{cases}
\phi(d(P,Q))\leq\phi(\tau), & \phi\text{ increasing},\\
\phi(d(P,Q))\geq\phi(\tau), & \phi\text{ decreasing}.
\end{cases}
\label{eq:threshold-equivalence}
\end{equation}
\end{lemma}

Strict monotonicity preserves or reverses order, so the converted threshold partitions the Jobs in the same way. The lemma applies to the Hellinger distance and its similarity forms, TVD and its complement, and absolute expectation deviation and observable fidelity. CCS belongs to a monotone class only above the chance baseline because clipping removes the order among lower values. DSR represents dominance among peaks and is not a transformation of task success, as Section~\ref{sec:dsr} establishes.

The coarse representation in Section~\ref{sec:primitives} provides an application to known answers. Hellinger distance and TVD are metrics on its $K+1$ bins~\cite{gilchrist2005,nielsen2010}. Applied to full histograms through this mapping, they are only pseudometrics whenever at least two outcomes outside $E$ are combined, because distinct histograms can then have zero separation.

\begin{proposition}[One accepted answer]
\label{prop:k1}
If $E$ contains one accepted answer with observed mass $p_E$, comparison of $(p_E,1-p_E)$ with $(1,0)$ gives
\begin{equation}
\TVD^{\mathrm{coarse}}=1-p_E,
\qquad
1-\TVD^{\mathrm{coarse}}=\HF^{\mathrm{coarse}}=p_E.
\label{eq:k1-coarse-relations}
\end{equation}
\end{proposition}

Substitution into the definitions of TVD and BC establishes these identities. Success rate, coarse TVD similarity, and coarse HF therefore provide alternative representations of the same task mass and should not be reported as independent evidence.

\begin{proposition}[Several accepted answers]
\label{prop:kmulti}
Let $K>1$, $p_i=\widehat P(e_i)$, and $w$ be the task weights. When $p_E>0$, let $r_i=p_i/p_E$ denote the conditional distribution of accepted answers. Then
\begin{equation}
\HF^{\mathrm{coarse}}=p_E\,\HF(r,w)\leq p_E.
\label{eq:multi-coarse-relations}
\end{equation}
Equality holds exactly when $r=w$. When $p_E=0$, coarse fidelity is zero. Coarse TVD similarity equals $p_E$ if and only if $p_i\leq w_i$ for every $i$.
\end{proposition}

Factoring $\sqrt{p_E}$ out of $\sum_i\sqrt{p_iw_i}$ gives Equation~\ref{eq:multi-coarse-relations}; the inner product inequality bounds $\HF(r,w)$ by one, with equality exactly at $r=w$. The TVD result follows from $1-\TVD^{\mathrm{coarse}}=\sum_i\min(p_i,w_i)$. Thus coarse fidelity combines accepted answer probability with agreement among accepted answers. A program change can improve this agreement without increasing task success.

The R1 task comparison coincides with the full R2 comparison when $P^{\star}=W_E$, because the ideal assigns zero mass outside $E$ and the same weights within it. An ideal Grover execution can retain probability outside $E$, so this condition need not hold. Subsection~\ref{sec:reference-evidence} examines the resulting differences in recorded executions.

\subsection{Task Success and Peak Dominance}
\label{sec:dsr}

DSR introduces another distinction within R1. It compares the average accepted peak $p_E/K$ with the largest competitor $p_{\mathrm{comp}}=\max_{x\notin E}\widehat P(x)$. For a valid count record with $0<K<2^m$, the clipped Michelson contrast~\cite{michelson1927} gives
\begin{equation}
\mathrm{DSR}=\max\left(0,\frac{p_E-Kp_{\mathrm{comp}}}{p_E+Kp_{\mathrm{comp}}}\right).
\label{eq:dsr-closed}
\end{equation}

For $p_{\mathrm{comp}}>0$, Equation~\ref{eq:dsr-closed} increases with the ratio of the average accepted peak to its strongest competitor above the clipping boundary. At or below that boundary it is zero; when $p_{\mathrm{comp}}=0$, it is one. DSR therefore represents a censored dominance ordering. This supports ordinal comparisons under a fixed answer set, outcome space, and sampling design~\cite{fentonBieman2015}, without identifying dominance with task success.

For example, with $m=10$ and $K=1$, distributions having $(p_E,p_{\mathrm{comp}})=(0.80,0.10)$ and $(0.08,0.01)$ both produce $7/9$. Equal DSR can therefore accompany task masses differing by a factor of ten. The score is unary for a fixed $E$, and taking the absolute difference of two scores yields only a pseudometric on distributions because distinct distributions can share a score. Changing $K$ alters the average accepted peak, while shot count and the number of competitors affect the observed maximum. These dependencies require justification before extending comparisons across tasks or sampling designs.

\subsection{Operations That Require Additional Assumptions}
\label{sec:operations}

Equivalent decisions for individual Jobs do not guarantee equivalent summaries over several Jobs. HF is the nonlinear transformation $\phi(H)=(1-H^2)^2$. Consider the constructed sets $A=\{0.10,0.10,0.90\}$ and $B=\{0.40,0.40,0.40\}$. Their summaries satisfy
\begin{equation}
\begin{aligned}
\operatorname{mean}(H_A)&=0.367<0.400=\operatorname{mean}(H_B),\\
\operatorname{mean}(\HF(A))&=0.665<0.706=\operatorname{mean}(\HF(B)).
\end{aligned}
\label{eq:aggregation-reversal}
\end{equation}
Smaller distance favors $A$, whereas larger fidelity favors $B$, so aggregation reverses the ordering. This is a mathematical counterexample rather than an observation from the execution corpus. By contrast, TVD similarity is an affine complement, so its arithmetic mean is exactly one minus the mean TVD and preserves the reversed direction. For an odd sample, the median is the central observation and transforms directly. The common median for an even sample averages two central observations and need not commute with a nonlinear transformation.

Processing measured outputs introduces a different requirement. TVD and Hellinger distance cannot increase when the same classical stochastic map, such as combining outcome bins, is applied to both distributions~\cite{csiszar1967}. The triangle inequality also bounds comparisons among three distributions on the same space. For TVD similarity it gives
\begin{equation}
\sigma_{\TVD}(P,R)
\geq
\sigma_{\TVD}(P,Q)+\sigma_{\TVD}(Q,R)-1.
\label{eq:tvd-similarity-triangle}
\end{equation}
Equation~\ref{eq:tvd-similarity-triangle} follows from TVD and becomes uninformative when its right side is negative. These properties concern classical distributions. Propagating error through quantum circuit stages requires suitable state or process measures and additional assumptions about their composition~\cite{gilchrist2005}.

For Estimator results, aggregation must follow the program's objective. If that objective is energy estimation for $H=\sum_i c_iO_i$, the signed error is $\Delta\langle H\rangle=\sum_i c_i\Delta\langle O_i\rangle$, and
\begin{equation}
\left|\Delta\langle H\rangle\right|
\leq
\sum_i|c_i|\left|\Delta\langle O_i\rangle\right|.
\label{eq:hamiltonian-error-bound}
\end{equation}
Equation~\ref{eq:hamiltonian-error-bound} explains why an unweighted average of observable fidelities does not generally represent energy error. Other objectives require their own observable weights or separate comparisons.

Scale claims impose a final restriction. A value in $[0,1]$ does not by itself support ratios, a common threshold across tasks, or an arithmetic summary without a defined estimand and weighting rule~\cite{fentonBieman2015,fenton1994}. A decision based on finite samples also requires an uncertainty rule. Table~\ref{tab:audit} consequently reports only the comparisons supported for each quantity instead of inferring them from its numerical range.

These results distinguish equivalent representations from changes in the evaluation itself. Converted thresholds preserve individual decisions, whereas nonlinear aggregation can change a group ordering and a different reference or attribute changes the question. The execution data next show which of these consequences appear in the recorded results.

\section{Evaluation Using QPU Execution Data}
\label{sec:evaluation}

% Internal review guide. This section supplies the empirical evidence for RQ3 by testing whether the distinctions established earlier appear in recorded executions and by showing how they affect the interpretation of DSR. The evidence is descriptive and does not benchmark providers.
The evaluation examines how the analytical distinctions affect recorded execution results. Corpus checks establish numerical consistency and identify whether aggregation changes a ranking. Grover cases then compare references and distinguish task success from DSR, before an Estimator execution shows how the same issue arises when reporting expectation estimates. The evidence concerns the interpretation of execution quantities, with no claim about relative device performance.

\subsection{Corpus Checks}
\label{sec:corpus-checks}

E1 checks whether stored Hellinger and TVD quantities satisfy Equations~\ref{eq:hellinger-fidelity} and~\ref{eq:tvd-similarity-definition}, including the reversal of rank. E2 searches for aggregation reversals between selected IBM and Rigetti groups. E3 counts records with task success available and full HF absent.

The Quantum Circuit Analysis and Runtime Development (QWARD) corpus contains 1,478 records of sampled outcomes from Grover, quantum Fourier transform, Bernstein--Vazirani, and teleportation executions, including four labeled as simulations~\cite{qwardGroverData2026}. The supplied \texttt{reproduce\_analysis.py} uses complete pairs for E1. E2 retains 559 IBM QPU records and 122 Rigetti Ankaa records from the designated execution groups, excluding the IonQ Forte record. It compares means and medians globally, by algorithm, by algorithm and qubit count, and within matched Grover configurations. General comparisons require five records per provider; matched configurations require one. These descriptive checks use overlapping groups and no significance test.

\begin{table*}[ht]
\centering
\caption{Reproducible corpus checks for transformation, aggregation, and reference availability.}
\label{tab:corpus-checks}
\renewcommand{\arraystretch}{1.1}
\begin{tabularx}{\textwidth}{@{}l>{\raggedright\arraybackslash}p{0.25\textwidth}
>{\raggedright\arraybackslash}p{0.32\textwidth}>{\raggedright\arraybackslash}X@{}}
\toprule
\textbf{Check} & \textbf{Rows or comparisons} & \textbf{Result} & \textbf{Interpretation} \\
\midrule
E1 Hellinger & 1,411 paired rows & maximum transform error $1.20\times10^{-6}$; $\rho=-0.999999999$ & stored values agree up to rounding \\
E1 TVD & 1,411 paired rows & maximum transform error $1.11\times10^{-16}$; $\rho=-1$ & exact complement to numerical precision \\
E2 aggregation & 31 provider comparisons for each class & zero mean inversions; zero median inversions & no empirical inversion in the specified groupings \\
E3 availability & all 1,478 rows & 67 with success rate present and full HF absent & an R1 score can remain when the R2 field is unavailable \\
\bottomrule
\end{tabularx}
\end{table*}

Table~\ref{tab:corpus-checks} establishes consistency of the stored representations. No aggregation reversal occurs in the 31 eligible groupings, so the mathematical counterexample in Equation~\ref{eq:aggregation-reversal} is not observed here. The 67 records with success rate present and full fidelity absent show different field availability; missingness alone does not establish why the R2 comparison was omitted.

\subsection{Reference Availability and Interpretation}
\label{sec:reference-evidence}

Availability limits which comparison can be reported, but even when both references are available they can encode different requirements. The supplied representation experiment illustrates the cost of one dense reference implementation; structured ideals may permit analytic or sparse computation. The Grover records provide the direct comparison of reference meaning.

Filtering to \texttt{algorithm = GROVER} yields 248 QPU records from IBM devices, Rigetti Ankaa, and IonQ Forte, with 246 complete pairs of coarse and full HF~\cite{qwardGroverData2026}. Their absolute difference reaches $0.231$. Among the 173 paired records with three or four qubits and 1,024 shots, applying the same numerical cutoff gives 13 disagreements at $\tau=0.7$ and 29 at $\tau=0.8$. In every disagreement, coarse fidelity fails while full fidelity passes. These illustrative cutoffs expose the effect of substituting a task target for the complete ideal. They are not equivalent acceptance criteria under Lemma~\ref{lem:threshold}, since the reference changes. For several accepted answers, coarse fidelity also includes agreement with the task weights by Proposition~\ref{prop:kmulti}; it cannot generally be identified with $p_E$.

\subsection{Grover with Multiple Marked States and DSR}
\label{sec:grover-multimarked}

The same Grover records show how a change of attribute affects interpretation. Table~\ref{tab:grover-multi} compares DSR with companion quantities for two four qubit configurations. Each entry is the median over optimization levels 0--3 on one backend. These summaries average the two central values and are not assumed to commute with nonlinear transformations, as discussed in Subsection~\ref{sec:operations}.

\begin{table*}[ht]
\centering
\caption{DSR and companion Grover measures after filtering the accompanying dataset to \texttt{algorithm = GROVER}. Values are medians over transpiler optimization levels 0--3 for each configuration and backend~\cite{qwardGroverData2026}.}
\label{tab:grover-multi}
\renewcommand{\arraystretch}{1.1}
\begin{tabular}{@{}lcccccc@{}}
\toprule
\textbf{Config.} & $K$ & $p_E$ & $\CCS$ & DSR & coarse HF & full HF \\
\midrule
\texttt{M4-2}, Torino & 2 & $0.517$ & $0.448$ & $0.675$ & $0.511$ & $0.734$ \\
\texttt{M4-4}, Fez & 4 & $0.403$ & $0.204$ & $0.000$ & $0.378$ & $0.378$ \\
\bottomrule
\end{tabular}
\end{table*}

For $K=2$, the complete ideal retains probability outside $E$, and the median full fidelity exceeds the median coarse fidelity. For $K=4$, one Grover iteration gives unit ideal task success, so $P^{\star}=W_E$ with uniform weights and the fidelities coincide. The median DSR is nevertheless zero while median task success is $0.403$. These are separate summaries of four executions. Three of those executions have zero DSR despite nonzero task success, illustrating the clipping and attribute distinction established by Equation~\ref{eq:dsr-closed}.

\subsection{Estimator Execution}
\label{sec:worked-estimator}

The distinction between agreement and task mass also arises for expectation estimates. A completed IBM Runtime EstimatorV2 execution evaluated a four qubit Greenberger--Horne--Zeilinger (GHZ) circuit on \texttt{ibm\_fez} against six Pauli observables with analytical R3 references~\cite{qwardEstimatorData2026}. Transpilation used optimization level 2 and produced instruction set architecture (ISA) depth 12. Table~\ref{tab:worked-estimator} selects three observables to compare expectation agreement with the affine transform of the estimated value.

\begin{table*}[ht]
\centering
\caption{Selected observables from the four qubit GHZ EstimatorV2 execution on \texttt{ibm\_fez}~\cite{qwardEstimatorData2026}.}
\label{tab:worked-estimator}
\renewcommand{\arraystretch}{1.1}
\begin{tabular}{@{}lcccc@{}}
\toprule
\textbf{Observable} & \textbf{Ideal} & $\widehat{\mathrm{ev}}\pm\mathrm{std}$ & \textbf{Obs.\ fidelity} & \textbf{Affine transform} \\
\midrule
$ZZZZ$ & $1$ & $0.9320\pm0.0086$ & $0.9660$ & $0.9660$ \\
$XZXZ$ & $0$ & $0.0005\pm0.0175$ & $0.9997$ & $0.5003$ \\
$YYYY$ & $1$ & $0.9175\pm0.0101$ & $0.9588$ & $0.9588$ \\
\bottomrule
\end{tabular}
\end{table*}

For $XZXZ$, Equation~\ref{eq:observable-fidelity} gives $0.9997$, indicating a small absolute deviation from the zero ideal. The stored record instead labels the affine value $(1+\widehat{\mathrm{ev}})/2=0.5003$ as \texttt{success\_prob\_4}. That label does not make it an agreement measure. Its interpretation as positive eigenspace mass requires the measurement assumptions in Section~\ref{sec:primitives}; the record omits mitigation metadata, so that interpretation cannot be verified here. The record also replaces the zero denominator in relative error with epsilon and reports approximately $5.37\times10^6$. The resulting number does not restore the undefined relative comparison.

These reporting choices expose a practical consequence of the classification. A developer consuming a field named success or error can draw a different conclusion from one based on the stated attribute and reference. Together with the Grover cases, this motivates making those choices explicit when selecting and reporting execution quantities.

\section{Discussion}
\label{sec:discussion}

% Internal review guide. This section synthesizes the answers to RQ1 through RQ3 as practical guidance. It should interpret the preceding analytical and execution evidence rather than introduce a new experiment or repeat the results.
The findings connect the mathematical properties of execution quantities to the requirements used to assess quantum software. Task success, distribution agreement, and peak dominance can each be useful, but they support different claims about the program. Selecting a quantity therefore requires stating the behavior to be evaluated and preserving that interpretation through the operations used to reach a decision.

\subsection{Selection of a Quantity}
\label{sec:selection}

Following GQM~\cite{basili1994}, the object is an individual execution, the purpose is evaluation, the quality focus is success for a stated task, and the viewpoint is the developer's. The context fixes the Primitive, available reference, shot budget, and mitigation procedure. Table~\ref{tab:gqm} derives questions and quantities from this goal using the classification in Table~\ref{tab:audit} and the operational conditions in Section~\ref{sec:analysis}.

\begin{table}[ht]
\centering
\caption{GQM questions and admissible quantities.}
\label{tab:gqm}
\renewcommand{\arraystretch}{1.12}
\begin{tabularx}{\columnwidth}{@{}>{\raggedright\arraybackslash}X
>{\raggedright\arraybackslash}X@{}}
\toprule
\textbf{Question} & \textbf{Quantity} \\
\midrule
How far is $\hat{P}$ from $P^{\star}$? & TVD or $H$ on $\Delta_{2^{m}}$ (R2) \\
How far is the coarse histogram from the weighted task target? & coarse TVD or coarse $H$ with justified $w$ (R1) \\
Did enough shots land in $E$? & $p_{E}$; $\CCS$ only with a justified null $Q_0$ \\
Do expected peaks dominate competitors? & DSR, an ordinal score of peak dominance (fixed $K$, $m$, and shot budget) \\
Is $P^{\star}$ available for the selected attribute? & If not, report R2 agreement as unavailable; use R1 only for a separately chosen answer attribute \\
How far is $\widehat{\mathrm{ev}}$ from the ideal? & $\lvert\Delta\mathrm{ev}\rvert$ or observable fidelity (R3) \\
Does the positive eigenspace define success? & if yes and the estimate is unmitigated, $p_+(O)$; otherwise do not label it success \\
Is uncertainty small enough to decide? & reported uncertainty and a decision interval; $z_0$ is only a standardized magnitude \\
Is the objective a Hamiltonian energy? & use its coefficients; report energy deviation and uncertainty \\
\bottomrule
\end{tabularx}
\end{table}

If a program must return any accepted answer, $p_E$ directly represents the observed frequency of that behavior. Requiring the accepted answers to follow specified weights adds a distribution requirement, as Proposition~\ref{prop:kmulti} establishes. DSR adds a different question about dominance over competitors. The Grover cases in Table~\ref{tab:grover-multi} demonstrate why these quantities can use the same counts without serving as interchangeable assessments of success. Their availability from sparse counts does not justify substituting them for an unavailable R2 comparison.

When assessing a program revision, the developer must also state how executions will be compared. A change from distance to its corresponding similarity requires the threshold conversion in Lemma~\ref{lem:threshold}. Aggregating several runs requires a defined summary and weighting because the nonlinear transformation in Equation~\ref{eq:aggregation-reversal} can change a ranking. For energy estimation, Equation~\ref{eq:hamiltonian-error-bound} connects the component deviations to the Hamiltonian objective. A decision under finite sampling additionally requires an uncertainty rule and may remain inconclusive when the available evidence cannot resolve the threshold.

The Estimator record in Table~\ref{tab:worked-estimator} makes these requirements relevant to software interfaces as well as analysis. A returned field should name its attribute and expose the reference and processing assumptions needed to interpret it. An undefined comparison should remain unavailable even when a numerical substitution produces a value. Combining quantities into a single score introduces another attribute and weighting scheme that require separate justification~\cite{fentonBieman2015}.

\subsection{Limitations}
\label{sec:threats}

% Internal review guide. This subsection limits the strength and scope of the conclusions, especially the execution evidence used for RQ3. It consolidates evidence, construct, generalization, literature, and researcher limitations without repeating the results.
The execution evidence supports descriptive conclusions for the recorded conditions. Backend, topology, compiler, execution date, and queue conditions were not controlled. Transformation checks establish numerical consistency, while aggregation uses existing records and overlapping groups selected by the authors. The absence of a reversal does not establish its prevalence, and the illustrative cutoffs do not account for sampling uncertainty. These observations therefore do not estimate a provider population effect.

The empirical cases cover a narrower domain than the classification. Detailed Sampler comparisons concern Grover, and the Estimator evidence contains one four qubit GHZ Job. The uniform null and task weights are study assumptions. The supplementary synthetic records test input processing rather than hardware accuracy. Agreement for these inputs and measured attributes does not establish correctness of the entire quantum program. Generalization to other workloads, measurement procedures, and tasks without suitable references requires further evidence; internal product measures and quantum advantage remain outside the scope.

The literature review was targeted rather than systematic, so differently framed or unindexed work could weaken the claim of a distinct integration. We also maintain the QWARD implementation that motivated the DSR case study, which creates a risk of confirmation bias. The supplied records and analysis programs permit inspection of the reported positive and negative findings, but independent replication and evaluation by practitioners are still needed to assess the classification and selection procedure more broadly.

\section{Conclusion}
\label{sec:conclusion}

% Internal review guide. Close the argument from the task requirement through the mathematical conditions to the observed consequences and their use in software development. Summarize RQ1 through RQ3 without repeating the selection checklist.
Success in a quantum execution depends on the behavior required by the program and the observations available to evaluate it. The classification in Section~\ref{sec:model} makes this relationship explicit by identifying the attribute, reference information, and permitted comparisons for each quantity. A shared numerical range does not make task success, distribution agreement, and peak dominance interchangeable.

The analytical results establish when a change of representation preserves the intended decision. Strictly monotone transformations preserve individual threshold decisions under the conditions of Lemma~\ref{lem:threshold}, while nonlinear aggregation can reverse a group ranking. Proposition~\ref{prop:kmulti} further shows that coarse HF combines accepted answer probability with agreement about their relative frequencies. DSR instead represents a censored ordering of peak dominance. These distinctions determine which mathematical operations answer the program's evaluation question.

The execution evidence in Section~\ref{sec:evaluation} illustrates the practical consequences. The corpus reproduces the defining transformations and contains no aggregation reversal in the selected groups. The Grover cases nevertheless produce different decisions under task and complete ideal references, while the Estimator record exposes a success label and a numerical error value that do not represent the intended agreement. The evidence supports these distinctions for the recorded conditions without establishing the correctness of the entire program.

For software development, the resulting contribution is a basis for choosing and interpreting execution quantities according to a stated requirement. The GQM procedure in Section~\ref{sec:selection} connects that requirement to the available reference and intended decision, including uncertainty and aggregation. Future work should evaluate this procedure with independent practitioners and controlled workloads to determine how reliably it supports quantum software assessment.

\section{Data and Code Availability}

The execution records and analysis programs are supplied with the article. The corpus checks use \texttt{DSR\_result.csv} and \texttt{reproduce\_analysis.py}~\cite{qwardGroverData2026}. The Grover examples apply the filter \texttt{algorithm = GROVER}; Table~\ref{tab:grover-multi} states the additional configuration, backend, and optimization level selections used for its two rows. The reference representation experiment uses \texttt{broad\_ideal\_experiment.py}, random seed 42, and the supplied JSON results. The Estimator example uses the completed IBM Runtime Job record in \texttt{estimator\_ibm\_result.json}~\cite{qwardEstimatorData2026}. These artifacts permit reproduction of the reported calculations without submitting a new QPU Job.

\section{Artificial Intelligence (AI) Tools and Services Acknowledgment}

The authors used AI assistance through different models for code development, quantum circuit analysis, statistical analysis, data processing, visualization, and manuscript revision. Codex supported this work, Writefull assisted with language and grammar improvements, and Elicit supported literature discovery. The authors retain responsibility for the study design, interpretation of results, and the accuracy and originality of the article.

\bibliographystyle{IEEEtran}
\bibliography{Bibliography}

\end{document}
